\scanpage{6}
car si \(x\) est un tel \unsure{élém.}, alors pour \(y\in C\) on a
\[
\unsure{\psi|\varphi}(xy)=\unsure{\psi|\varphi}(x)y=0,
\]
d'où (comme \(xy\in C\), \unsure{alors|déjà} \(\unsure{\psi|\varphi}(xy)=xy\)) \(xy=0\), et ceci pour tout \(y\in C\), d'où \(x=0\). -- \uline{Donc \(\psi\) est un isomorphisme de \(\Gamma\) sur \(\Gamma_1\)}, \deleted{(\unsure{et le} \(\varphi\) \unsure{est la réciproque})}\inserted{\deleted{\unsure{et}} dont \(\varphi\)} \uline{est l'inverse}.

\unsure{Il n'y a plus} qu'à prouver que \(\psi\varphi=1\), soit donc \(u\in\Gamma_1\)
\[
\varphi u=v\in C''\quad\text{est donné par}\quad\langle v,f\rangle=\uline{\langle u,f\rangle}\quad(f\in C')
\]
\deleted{\unsure{d'où|alors}} \(\psi v\) est donné par
\[
\langle\psi v,f\rangle=\langle v,f_1\rangle\quad(f\in R_*)
\]
i.e.
\[
\langle\psi v,f\rangle=\langle u,f_1\rangle\quad(f\in R_*),
\]
mais on a déjà dit que \(\langle u,f_1\rangle=\langle u,f\rangle\), d'où
\[
\langle\psi v,f\rangle=\langle u,f\rangle,\quad\text{i.e. }\psi v=u.\qquad\text{CQFD.}
\]

\mspage{7}
